\documentclass[a4paper,11pt,dvips]{article} 
\usepackage[T1]{fontenc}
\usepackage[utf8]{inputenc}
\usepackage{lmodern}
\usepackage{geometry}
\geometry{hmargin=1.2cm,vmargin=1.2cm}
\usepackage{graphicx}
\usepackage{pstricks,pst-plot,pst-tree}
%\usepackage{pst-eucl} 
\usepackage{framed}
\usepackage{amsmath,bm}
\usepackage{amssymb}
\usepackage{fancyhdr}
\usepackage{fancybox}
\usepackage{multicol}
\usepackage{xcolor}
\usepackage{epsfig}
\usepackage{pifont}
\usepackage[framed]{ntheorem}
\usepackage[frenchb]{babel}
\usepackage{tabularx}
\usepackage{pstricks,pst-plot,pst-eucl}
\usepackage{pst-tree}
\usepackage{mathtools}
\usepackage{amssymb}
\usepackage{babel}



\begin{document}
\pagestyle{empty}
\textbf{\begin{center}
\Huge{Devoir surveillé V}
\end{center}}
\begin{center}
\textit{Dans tout le devoir, on donnera les résultats sous forme de fraction irréductible.}
\end{center}


\section{Calcul de probabilités}
   On tire au hasard une carte dans un jeu de 32 cartes.
   \begin{enumerate}
      \item Quelle est la probabilité de tirer un trèfle ?
      \item Quelle est la probabilité de tirer une carte noire     ?
      \item Quelle est la probabilité de ne pas tirer un carreau ?
      \item Quelle est la probabilité de tirer une figure (roi, dame ou valet) ?
      \item Quelle est la probabilité de tirer un as ?
      \item Quelle est la probabilité de ne pas tirer un valet noir ? \\
   \end{enumerate}

\section{Loi de probabilité}
On lance deux dés cubiques équilibrés numérotés de 1 à 6.    On note alors le plus grand des deux numéros sortis.
   \begin{enumerate}
      \item Utiliser un tableau à double entrée pour modéliser la situation.
      \item Quel est l'univers $\Omega$ de toutes les issues possibles ?
      \item Établir la loi de probabilité de l'expérience.
   \end{enumerate}

\section{Arbre pondéré}

Une urne contient 5 boules indiscernables au toucher : deux bleues et trois rouges.
On dispose également de deux sacs contenant des jetons : l'un est bleu et contient un jeton 
bleu et trois jetons rouges, l’autre est rouge et contient deux jetons bleus et deux jetons rouge. \\

On note l'événement:
\begin{itemize}
\item B :"Tirer une boule bleu"
\item R :"Tirer une boule rouge"
\item b :"Tirer un jeton bleu"
\item r :"Tirer un jeton rouge"\\
\end{itemize}

On extrait une boule de l'urne, puis on tire un jeton dans le sac qui est de la même couleur que la boule tirée.

\begin{enumerate}
\item Combien y a-t-il d’issues possibles ?
\item A l'aide d’un arbre pondéré, quel est la probabilité de chacune de ses issues.
\item Quel est la probabilité de l'événement A : « la boule et le jeton extraits sont de la même
couleur »
\end{enumerate}


\end{document}